% Version 2 worked exam, adapted from G11_T1_L3_C8C9C10_exam_solution.tex.
% C8/C9 content: the year-local C8 and C9 classroom activities.
% C10 content: source exam. Presentation: existing solution-document macros.
\makeatletter
\def\input@path{{./}{../}{../../}{../../../}{../../../../}{preamble/}{../preamble/}{../../preamble/}{../../../preamble/}{../../../../preamble/}}
\makeatother
\input{preamble_exams.tex}
\graphicspath{{./}{../}{../../}{../../../}{../../../../}{preamble/}{../preamble/}{../../preamble/}{../../../preamble/}{../../../../preamble/}}
\linespread{1.0}
\setlength{\parskip}{5pt}
\usepackage{needspace}
\title{Grade 11 --- Learning Evidence 3\\C8, C9 and C10 --- Exam Solutions --- Version 2\\\large 1st Term \quad 2026--2027}
\author{}
\date{}
\begin{document}
\maketitle
\setcounter{tocdepth}{1}
\setcounter{secnumdepth}{2}
\phantomsection
\label{toc}
\tableofcontents

\section*{Learning objective}
Apply continuous probability models and sampling distributions to investigate economic and financial questions, justify statistical methods, and communicate evidence-based conclusions.

\section*{Criteria assessment}
Partial credit may be awarded when the correct procedure is clearly demonstrated, even if arithmetic errors are made. No points will be awarded for answers without a demonstrated procedure.

A minimum of 4 points is required to pass each criterion. A 1-point penalty will be applied to each criterion if the student's name is not written on the assessment or if the assessment is submitted late.


\begin{center}
	\begin{tabular}{|c|c|}
		\hline
		\textbf{\textnumero} & \textbf{C8} \\ \hline
		1--5 & 5 \\ \hline
		6--10 & 5 \\ \hline
		\textbf{Total} & \textbf{10} \\ \hline
	\end{tabular}
	\hspace{1cm}
	\begin{tabular}{|c|c|}
		\hline
		\textbf{\textnumero} & \textbf{C9} \\ \hline
		11--13 & 6 \\ \hline
		14--16 & 6 \\ \hline
		\textbf{Total} & \textbf{12} \\ \hline
	\end{tabular}
\end{center}


The worked solutions show the reasoning using the supplied normal and Student's $t$ reference sheets; no interpolation, CDF calculator, or inverse-CDF calculator is needed.
C8 is worth \textbf{10 points}, C9 \textbf{12 points}, and C10 \textbf{6 points}: \textbf{28 points total}.
At least \textbf{4 points per criterion} are required to pass that criterion.
All calculations within one problem share its stated allocation, including both interval endpoints.

\section*{Evaluated criteria}
\begin{enumerate}[label=\textbf{C\arabic*},start=8,leftmargin=3em]
\item Uses Student's t-distribution tables, degrees of freedom, and one-tailed or two-tailed areas to determine critical values, central probabilities, and tail probabilities or probability ranges.
\item Calculates and compares critical values and probability areas from the standard normal and Student's t-distributions, selecting and justifying the appropriate distribution according to whether the population standard deviation is known or estimated from the sample.
\item Integrates continuous probability modelling, simulation, sampling analysis, and portfolio-risk measures to investigate an economic or financial question, justifying the selected methods and communicating results, assumptions, limitations, and evidence-based conclusions.
\end{enumerate}

\section*{Models and table conventions}
Problems 1--10 use the supplied signed-deviation model $X=m+aT$, $T\sim t_\nu$.
Problems 11--16 compare
\[
 X_T=m+aT,\quad T\sim t_\nu,\qquad X_N=m+aZ,\quad Z\sim N(0,1),
 \qquad \boxed{\nu=n-1}.
\]
Here $n$ is the number of sample observations, $m$ is the location (and mean for $\nu>1$), and $a>0$ is the scale.
Use the stated \emph{approximate standard deviation} directly as $a$, without adjustment.
This is a teaching approximation:
\[
 \operatorname{SD}(X_T)=a\sqrt{\frac{\nu}{\nu-2}}\quad(\nu>2),
 \qquad \operatorname{SD}(X_N)=a.
\]
The two models share a scale, not an actual variance. The approximation is coarser for small $\nu$ and improves as $\nu$ increases.
Negative deviations mean below the contextual reference (early or under budget); positive deviations mean above it (late or over budget).
Inverse intervals enclose modelled deviations. Standardize and back-transform using
\[
 \boxed{c=\frac{x-m}{a}},\qquad \boxed{x=m+ac}.
\]
The year-local sheets are in \texttt{book2026-2027/fundamentals/normal-table.tex} (and its PDF).
In the $t$ table, row $\nu$ and one-tail column $\alpha$ give positive $t^*$ satisfying
\[
 P(T>t^*)=\alpha,\qquad P(T<-t^*)=\alpha,\qquad
 P(-t^*<T<t^*)=1-2\alpha.
\]
Obtain negative cutoffs by symmetry. The printed columns are
$0.30,0.20,0.10,0.05,0.025,0.01,0.005$.
Direct bounds recover printed three-decimal entries exactly through standardization.
Report C8 probabilities as implied by these entries, and contextual cutoffs to two decimals.

The normal table gives four-decimal cumulative areas $\Phi(z)=P(Z\le z)$ for positive $z$.
Use $\Phi(-z)=1-\Phi(z)$. For direct normal probabilities, round each standardized boundary to two decimals \emph{before} lookup.
For inverse normal problems, select the closest printed CDF entry.
At targets $0.0500$ and $0.9500$, magnitudes $1.64$ and $1.65$ tie; choose $|z|=1.64$, as in the C9 activity.
Use the CDF table for inverse normal lookups.

For \textbf{every C9 problem}, calculate both results and
\[
 E_A=|r_T-r_N|,\qquad E_{Pt}=\frac{100E_A}{|r_T|},\qquad
 E_{Pn}=\frac{100E_A}{|r_N|}.
\]
For direct problems, $r_T,r_N$ are probabilities; for inverse problems, they are contextual cutoffs after back-transformation.
Compare inverse interval endpoints separately.
A zero denominator makes that percentage undefined; the data below avoid zero denominators.
These measures compare model results and do not establish error against observed truth.
Report C9 probabilities and their $E_A$ to four decimals; contextual cutoffs and their $E_A$ to two decimals with units; and percentage comparisons to two decimals.
Use unrounded back-transformed cutoffs for inverse comparisons, rounding only final results.

All results are table-based approximations, not exact distribution values: printed critical values and CDF entries are rounded.
Equalities in the worked arithmetic use selected table values as working values.
Both models are continuous, so endpoints have probability zero.
The formally stated criteria above are retained; the supplied activity models govern the C8/C9 calculations here.


\needspace{8\baselineskip}
\sectionScored{Warehouse energy deviations}{C8-\sml{1}}{sec:q1}
{
\textbf{C8 -- Direct interval above the mean.} Over $n=18$ days, electricity-use deviation from forecast, in kWh, has mean $m=2.4$ and approximate standard deviation $a=7$. Use $X=2.4+7T$, $T\sim t_{n-1}$.
\QuestionsTasks{
{8}/{sol:q1}/{\textbf{Problem 1.} Find $P(8.441<X<14.58)$. \score{1}}
}
\needspace{5\baselineskip}
\phantomsection
\label{sol:q1}
\subsection*{C8 -- Problem 1: Direct interval above the mean solution \score{1}}
Here $\nu=18-1=17$. Both bounds exceed $m=2.4$:
\[
 c_L=\frac{8.441-2.4}{7}=0.863,\qquad c_U=\frac{14.58-2.4}{7}=1.740.
\]
Row $17$, columns $\alpha=0.20$ and $0.05$, gives $0.863$ and $1.740$.
Subtract the nested right tails:
\[
 P(8.441<X<14.58)\approx0.20-0.05=\boxed{0.15}.
\]
About $15\%$ of days use between $8.44$ and $14.58$ kWh above forecast.
\secInicio[sec:q1]
}

\needspace{8\baselineskip}
\sectionScored{Retail revenue alerts}{C8-\sml{1}}{sec:q2}
{
\textbf{C8 -- Direct unequal combined tails.} Over $n=8$ weeks, revenue deviation from plan, in thousands of dollars, has mean $m=1.6$ and approximate standard deviation $a=5$. Use $X=1.6+5T$, $T\sim t_{n-1}$.
\QuestionsTasks{
{8}/{sol:q2}/{\textbf{Problem 2.} Find the probability of an alert: $X<-7.875$ or $X>8.675$. \score{1}}
}
\needspace{5\baselineskip}
\phantomsection
\label{sol:q2}
\subsection*{C8 -- Problem 2: Direct unequal combined tails solution \score{1}}
Here $\nu=8-1=7$ and
\[
 c_L=\frac{-7.875-1.6}{5}=-1.895,\qquad c_U=\frac{8.675-1.6}{5}=1.415.
\]
Row $7$, columns $\alpha=0.05$ and $0.10$, gives $1.895$ and $1.415$.
Add the disjoint tails:
\[
 P(X<-7.875\text{ or }X>8.675)\approx0.05+0.10=\boxed{0.15}.
\]
About $15\%$ of weeks trigger an alert: revenue is more than $7.88$ thousand dollars below plan or $8.68$ thousand dollars above it.
\secInicio[sec:q2]
}

\needspace{8\baselineskip}
\sectionScored{Courier arrival deviations}{C8-\sml{1}}{sec:q3}
{
\textbf{C8 -- Direct left tail.} A sample of $n=11$ courier trips has arrival-time deviation from schedule, in minutes, with mean $m=1.2$ and approximate standard deviation $a=3.5$. Use $X=1.2+3.5T$, $T\sim t_{n-1}$.
\QuestionsTasks{
{8}/{sol:q3}/{\textbf{Problem 3.} Find $P(X<-5.142)$. \score{1}}
}
\needspace{5\baselineskip}
\phantomsection
\label{sol:q3}
\subsection*{C8 -- Problem 3: Direct left tail solution \score{1}}
Here $\nu=11-1=10$. Standardize:
\[
 c=\frac{-5.142-1.2}{3.5}=-1.812.
\]
Row $10$, column $\alpha=0.05$, gives $t^*=1.812$.
By symmetry,
\[
 P(X<-5.142)=P(T<-1.812)\approx\boxed{0.05}.
\]
About $5\%$ of trips arrive more than $5.14$ minutes early.
\secInicio[sec:q3]
}

\needspace{8\baselineskip}
\sectionScored{Freight cost deviations}{C8-\sml{1}}{sec:q4}
{
\textbf{C8 -- Direct right tail.} A sample of $n=14$ consignments has cost deviation from budget, in dollars, with mean $m=-5$ and approximate standard deviation $a=9$. Use $X=-5+9T$, $T\sim t_{n-1}$.
\QuestionsTasks{
{8}/{sol:q4}/{\textbf{Problem 4.} Find $P(X>-12.83)$. \score{1}}
}
\needspace{5\baselineskip}
\phantomsection
\label{sol:q4}
\subsection*{C8 -- Problem 4: Direct right tail solution \score{1}}
Here $\nu=14-1=13$ and
\[
 c=\frac{-12.83-(-5)}{9}=-0.870.
\]
Row $13$, column $\alpha=0.20$, gives $t^*=0.870$.
Symmetry gives $P(T<-0.870)\approx0.20$, so
\[
 P(X>-12.83)=1-P(T<-0.870)\approx1-0.20=\boxed{0.80}.
\]
About $80\%$ have costs above the threshold of approximately $\$12.83$ below budget.
The cutoff is below $m=-5$, so the right-tail event includes more than half the symmetric model.
\secInicio[sec:q4]
}

\needspace{8\baselineskip}
\sectionScored{Crop shipment deviations}{C8-\sml{1}}{sec:q5}
{
\textbf{C8 -- Direct interval below the mean.} A sample of $n=9$ crop shipments has mass deviation from forecast, in tonnes, with mean $m=-1.2$ and approximate standard deviation $a=3.6$. Use $X=-1.2+3.6T$, $T\sim t_{n-1}$.
\QuestionsTasks{
{8}/{sol:q5}/{\textbf{Problem 5.} Find $P(-9.5016<X<-4.4004)$. \score{1}}
}
\needspace{5\baselineskip}
\phantomsection
\label{sol:q5}
\subsection*{C8 -- Problem 5: Direct interval below the mean solution \score{1}}
Here $\nu=9-1=8$. Both bounds lie below $m=-1.2$:
\[
 c_L=\frac{-9.5016-(-1.2)}{3.6}=-2.306,\qquad c_U=\frac{-4.4004-(-1.2)}{3.6}=-0.889.
\]
Row $8$, columns $\alpha=0.025$ and $0.20$, gives $2.306$ and $0.889$.
By symmetry the cumulative areas at the negative bounds are $0.025$ and $0.20$:
\[
 P(-9.5016<X<-4.4004)\approx0.20-0.025=\boxed{0.175}.
\]
About $17.5\%$ of shipments fall short of forecast by between $4.40$ and $9.50$ tonnes.
\secInicio[sec:q5]
}

\needspace{8\baselineskip}
\sectionScored{Package mass alerts}{C8-\sml{1}}{sec:q6}
{
\textbf{C8 -- Inverse equal-tail alert limits.} A sample of $n=10$ packages has mass deviation from target, in grams, with mean $m=-0.5$ and approximate standard deviation $a=1.4$. Use $X=-0.5+1.4T$, $T\sim t_{n-1}$.
\QuestionsTasks{
{8}/{sol:q6}/{\textbf{Problem 6.} Find two alert limits leaving a combined outside probability of $0.05$, split equally. \score{1}}
}
\needspace{5\baselineskip}
\phantomsection
\label{sol:q6}
\subsection*{C8 -- Problem 6: Inverse equal-tail alert limits solution \score{1}}
Here $\nu=10-1=9$. Each tail has probability $0.05/2=0.025$.
Row $9$, column $\alpha=0.025$, gives $t^*=2.262$, so $(c_L,c_U)=(-2.262,2.262)$:
\[
 \begin{aligned}
 L&=-0.5-1.4(2.262)=-3.6668\approx\boxed{-3.67\text{ g}},\\
 U&=-0.5+1.4(2.262)=2.6668\approx\boxed{2.67\text{ g}}.
 \end{aligned}
\]
About $2.5\%$ of packages fall below $L$ and $2.5\%$ above $U$.
An alert means more than $3.67$ g under target or $2.67$ g over target.
\secInicio[sec:q6]
}

\needspace{8\baselineskip}
\sectionScored{Operating cost cutoff}{C8-\sml{1}}{sec:q7}
{
\textbf{C8 -- Inverse lower-tail cutoff.} Over $n=16$ production days, operating-cost deviation from budget, in dollars, has mean $m=-3.5$ and approximate standard deviation $a=6$. Use $X=-3.5+6T$, $T\sim t_{n-1}$.
\QuestionsTasks{
{8}/{sol:q7}/{\textbf{Problem 7.} Find $x$ such that $P(X<x)=0.10$. \score{1}}
}
\needspace{5\baselineskip}
\phantomsection
\label{sol:q7}
\subsection*{C8 -- Problem 7: Inverse lower-tail cutoff solution \score{1}}
Here $\nu=16-1=15$. The lower-tail target is $0.10$.
Row $15$, column $\alpha=0.10$, gives $t^*=1.341$.
By symmetry $c=-1.341$; back-transform:
\[
 x=-3.5+6(-1.341)=-11.546\approx\boxed{-11.55\text{ dollars}}.
\]
About $10\%$ of days have costs more than $\$11.55$ below budget.
\secInicio[sec:q7]
}

\needspace{8\baselineskip}
\sectionScored{Dispatch time alerts}{C8-\sml{1}}{sec:q8}
{
\textbf{C8 -- Inverse unequal-tail alert limits.} A sample of $n=17$ dispatches has departure-time deviation from schedule, in minutes, with mean $m=-2.5$ and approximate standard deviation $a=7$. Use $X=-2.5+7T$, $T\sim t_{n-1}$.
\QuestionsTasks{
{8}/{sol:q8}/{\textbf{Problem 8.} Find $L,U$ with $P(X<L)=0.01$ and $P(X>U)=0.10$. \score{1}}
}
\needspace{5\baselineskip}
\phantomsection
\label{sol:q8}
\subsection*{C8 -- Problem 8: Inverse unequal-tail alert limits solution \score{1}}
Here $\nu=17-1=16$. Use the tails separately.
Row $16$, column $\alpha=0.01$, gives $2.583$, so $c_L=-2.583$ by symmetry.
Column $\alpha=0.10$ gives $c_U=1.337$:
\[
 \begin{aligned}
 L&=-2.5+7(-2.583)=-20.581\approx\boxed{-20.58\text{ minutes}},\\
 U&=-2.5+7(1.337)=6.859\approx\boxed{6.86\text{ minutes}}.
 \end{aligned}
\]
About $1\%$ depart more than $20.58$ minutes early, and $10\%$ more than $6.86$ minutes late.
The combined alert probability is $0.01+0.10=0.11$; $0.89$ lies between the limits.
\secInicio[sec:q8]
}

\needspace{8\baselineskip}
\sectionScored{Loading time cutoff}{C8-\sml{1}}{sec:q9}
{
\textbf{C8 -- Inverse upper-tail cutoff.} A sample of $n=22$ loading jobs has duration deviation from target, in minutes, with mean $m=2.5$ and approximate standard deviation $a=8$. Use $X=2.5+8T$, $T\sim t_{n-1}$.
\QuestionsTasks{
{8}/{sol:q9}/{\textbf{Problem 9.} Find $x$ such that $P(X>x)=0.025$. \score{1}}
}
\needspace{5\baselineskip}
\phantomsection
\label{sol:q9}
\subsection*{C8 -- Problem 9: Inverse upper-tail cutoff solution \score{1}}
Here $\nu=22-1=21$. Row $21$, column $\alpha=0.025$, gives $c=2.080$:
\[
 P(T>2.080)\approx0.025,\qquad
 x=2.5+8(2.080)=19.140\approx\boxed{19.14\text{ minutes}}.
\]
About $2.5\%$ of jobs take more than $19.14$ minutes longer than target.
\secInicio[sec:q9]
}

\needspace{8\baselineskip}
\sectionScored{Demand deviation interval}{C8-\sml{1}}{sec:q10}
{
\textbf{C8 -- Inverse central interval.} Over $n=13$ weeks, demand deviation from forecast, in tonnes, has mean $m=5$ and approximate standard deviation $a=10$. Use $X=5+10T$, $T\sim t_{n-1}$.
\QuestionsTasks{
{8}/{sol:q10}/{\textbf{Problem 10.} Find $L<U$ enclosing the central $90\%$ of deviations with equal excluded tails. \score{1}}
}
\needspace{5\baselineskip}
\phantomsection
\label{sol:q10}
\subsection*{C8 -- Problem 10: Inverse central interval solution \score{1}}
Here $\nu=13-1=12$. Each tail is $(1-0.90)/2=0.05$.
Row $12$, column $\alpha=0.05$, gives $t^*=1.782$, so $c_L=-1.782$, $c_U=1.782$:
\[
 \begin{aligned}
 L&=5-10(1.782)=-12.820\approx\boxed{-12.82\text{ tonnes}},\\
 U&=5+10(1.782)=22.820\approx\boxed{22.82\text{ tonnes}}.
 \end{aligned}
\]
About $90\%$ lie between $12.82$ tonnes below forecast and $22.82$ tonnes above it, with $5\%$ in each excluded tail.
\secInicio[sec:q10]
}

\needspace{8\baselineskip}
\sectionScored{Delivery schedule deviations}{C9-\sml{2}}{sec:q11}
{
\textbf{C9 -- Direct right tail.} A sample of $n=13$ deliveries has arrival-time deviation from schedule, in minutes, with mean $m=3.2$ and approximate standard deviation $a=4.5$. Compare $X_T=3.2+4.5T$, $T\sim t_{n-1}$, with $X_N=3.2+4.5Z$.
\QuestionsTasks{
{9}/{sol:q11}/{\textbf{Problem 11.} Find $P(X>-0.7285)$. Find both model results, calculate $E_A$, $E_{Pt}$, and $E_{Pn}$. \score{2}}
}
\needspace{5\baselineskip}
\phantomsection
\label{sol:q11}
\subsection*{C9 -- Problem 11: Direct right tail solution \score{2}}
Here $\nu=13-1=12$; $c=(-0.7285-3.2)/4.5=-0.873$.
Row $12$, column $\alpha=0.20$, gives $t^*=0.873$:
\[
 p_T=1-P(T<-0.873)\approx1-0.2000=\boxed{0.8000}.
\]
For normal, round to $z=-0.87$.
Row $0.8$, column $0.07$, gives $\Phi(0.87)=0.8078$:
\[
 \Phi(-0.87)=1-0.8078=0.1922,\qquad p_N=1-0.1922=\boxed{0.8078}.
\]
\[
 \begin{aligned}
 E_A&=|0.8000-0.8078|=\boxed{0.0078},\\
 E_{Pt}&=\frac{100(0.0078)}{|0.8000|}\approx\boxed{0.98\%},\\
 E_{Pn}&=\frac{100(0.0078)}{|0.8078|}\approx\boxed{0.97\%}.
 \end{aligned}
\]
The event includes arrivals less than about $0.73$ minutes early and all late arrivals.
Normal assigns $80.78\%$, which is $0.78$ percentage points higher than $t$'s $80.00\%$.
Both exceed $0.5$ because the cutoff is below $m=3.2$; $t$ does not give the larger probability for this event.
\secInicio[sec:q11]
}

\needspace{8\baselineskip}
\sectionScored{Fuel forecast deviations}{C9-\sml{2}}{sec:q12}
{
\textbf{C9 -- Direct interval crossing the mean.} Over $n=14$ operating days, fuel-use deviation from forecast, in litres, has mean $m=1.2$ and approximate standard deviation $a=7$. Compare $X_T=1.2+7T$, $T\sim t_{n-1}$, with $X_N=1.2+7Z$.
\QuestionsTasks{
{9}/{sol:q12}/{\textbf{Problem 12.} Find $P(-8.25<X<16.32)$. Find both model results, calculate $E_A$, $E_{Pt}$, and $E_{Pn}$. \score{2}}
}
\needspace{5\baselineskip}
\phantomsection
\label{sol:q12}
\subsection*{C9 -- Problem 12: Direct interval crossing the mean solution \score{2}}
Here $\nu=14-1=13$. Standardize:
\[
 c_L=\frac{-8.25-1.2}{7}=-1.350,\qquad c_U=\frac{16.32-1.2}{7}=2.160.
\]
Row $13$, columns $\alpha=0.10$ and $0.025$, gives $1.350$ and $2.160$.
The interval crosses $m=1.2$ and excludes unequal tails:
\[
 p_T\approx1-0.1000-0.0250=\boxed{0.8750}.
\]
For normal, round to $z_L=-1.35$ and $z_U=2.16$.
Row $1.3$, column $0.05$, gives $\Phi(1.35)=0.9115$; row $2.1$, column $0.06$, gives $\Phi(2.16)=0.9846$:
\[
 p_N=0.9846-(1-0.9115)=0.9846-0.0885=\boxed{0.8961}.
\]
\[
 \begin{aligned}
 E_A&=|0.8750-0.8961|=\boxed{0.0211},\\
 E_{Pt}&=\frac{100(0.0211)}{|0.8750|}\approx\boxed{2.41\%},\\
 E_{Pn}&=\frac{100(0.0211)}{|0.8961|}\approx\boxed{2.35\%}.
 \end{aligned}
\]
Fuel use lies from $8.25$ litres below forecast to $16.32$ litres above it with probability about $87.50\%$ under $t$ or $89.61\%$ under normal.
Normal assigns $2.11$ percentage points more to this interval.
\secInicio[sec:q12]
}

\needspace{8\baselineskip}
\sectionScored{Sales revenue shortfalls}{C9-\sml{2}}{sec:q13}
{
\textbf{C9 -- Direct left tail.} Over $n=11$ weeks, sales-revenue deviation from target, in thousands of dollars, has mean $m=-2.5$ and approximate standard deviation $a=6$. Compare $X_T=-2.5+6T$, $T\sim t_{n-1}$, with $X_N=-2.5+6Z$.
\QuestionsTasks{
{9}/{sol:q13}/{\textbf{Problem 13.} Find $P(X<-13.372)$. Find both model results, calculate $E_A$, $E_{Pt}$, and $E_{Pn}$. \score{2}}
}
\needspace{5\baselineskip}
\phantomsection
\label{sol:q13}
\subsection*{C9 -- Problem 13: Direct left tail solution \score{2}}
Here $\nu=11-1=10$; $c=(-13.372-(-2.5))/6=-1.812$.
Row $10$, column $\alpha=0.05$, gives $t^*=1.812$; symmetry gives
\[
 p_T=P(T<-1.812)\approx\boxed{0.0500}.
\]
For normal, round to $z=-1.81$. Row $1.8$, column $0.01$, gives $\Phi(1.81)=0.9649$:
\[
 p_N=1-0.9649=\boxed{0.0351}.
\]
\[
 \begin{aligned}
 E_A&=|0.0500-0.0351|=\boxed{0.0149},\\
 E_{Pt}&=\frac{100(0.0149)}{|0.0500|}\approx\boxed{29.80\%},\\
 E_{Pn}&=\frac{100(0.0149)}{|0.0351|}\approx\boxed{42.45\%}.
 \end{aligned}
\]
A shortfall greater than $13.37$ thousand dollars has probability about $5.00\%$ under $t$, compared with $3.51\%$ under normal, a difference of $1.49$ percentage points.
\secInicio[sec:q13]
}

\needspace{8\baselineskip}
\sectionScored{Cargo mass deviation interval}{C9-\sml{2}}{sec:q14}
{
\textbf{C9 -- Inverse central interval.} A sample of $n=11$ cargo loads has mass deviation from target, in kg, with mean $m=1.8$ and approximate standard deviation $a=5$. Compare $X_T=1.8+5T$, $T\sim t_{n-1}$, with $X_N=1.8+5Z$.
\QuestionsTasks{
{9}/{sol:q14}/{\textbf{Problem 14.} Find the central $95\%$ interval with equal excluded tails under each model. Calculate $E_A$, $E_{Pt}$, and $E_{Pn}$ separately for each endpoint. \score{2}}
}
\needspace{5\baselineskip}
\phantomsection
\label{sol:q14}
\subsection*{C9 -- Problem 14: Inverse central interval solution \score{2}}
Here $\nu=11-1=10$. Each tail is $(1-0.95)/2=0.025$, so cumulative targets are $0.0250$ and $0.9750$.
Standardize unknown limits with $c_L=(L-1.8)/5$ and $c_U=(U-1.8)/5$.
Row $10$, column $\alpha=0.025$, gives $t^*=2.228$, hence $(c_{T,L},c_{T,U})=(-2.228,2.228)$.
For normal, row $1.9$, column $0.06$, gives $\Phi(1.96)=0.9750$ exactly at printed precision; symmetry gives $\Phi(-1.96)=0.0250$.
Thus $(c_{N,L},c_{N,U})=(-1.96,1.96)$.

Back-transform, retaining unrounded values:
\[
 \begin{aligned}
 L_T&=1.8-5(2.228)=-9.340\approx\boxed{-9.34\text{ kg}},\\
 U_T&=1.8+5(2.228)=12.940\approx\boxed{12.94\text{ kg}},\\
 L_N&=1.8-5(1.96)=-8.00=\boxed{-8.00\text{ kg}},\\
 U_N&=1.8+5(1.96)=11.60=\boxed{11.60\text{ kg}}.
 \end{aligned}
\]
\textbf{Lower endpoint comparison:}
\[
 \begin{aligned}
 E_{A,L}&=|-9.340-(-8.00)|=1.340\approx\boxed{1.34\text{ kg}},\\
 E_{Pt,L}&=\frac{100(1.340)}{|-9.340|}\approx\boxed{14.35\%},\\
 E_{Pn,L}&=\frac{100(1.340)}{|-8.00|}\approx\boxed{16.75\%}.
 \end{aligned}
\]
\textbf{Upper endpoint comparison:}
\[
 \begin{aligned}
 E_{A,U}&=|12.940-11.60|=1.340\approx\boxed{1.34\text{ kg}},\\
 E_{Pt,U}&=\frac{100(1.340)}{|12.940|}\approx\boxed{10.36\%},\\
 E_{Pn,U}&=\frac{100(1.340)}{|11.60|}\approx\boxed{11.55\%}.
 \end{aligned}
\]
The central intervals are approximately $(-9.34,12.94)$ kg under $t$ and $(-8.00,11.60)$ kg under normal, with $2.5\%$ in each excluded tail.
Negative limits mean below target mass; positive limits mean above it.
The normal interval extends about $1.34$ kg less far from $m$ at each end.
Percentage comparisons differ because the endpoint magnitudes differ.
\secInicio[sec:q14]
}

\needspace{8\baselineskip}
\sectionScored{Production cost savings}{C9-\sml{2}}{sec:q15}
{
\textbf{C9 -- Inverse lower-tail cutoff.} A sample of $n=9$ batches has cost deviation from budget, in dollars, with mean $m=-1.3$ and approximate standard deviation $a=3$. Compare $X_T=-1.3+3T$, $T\sim t_{n-1}$, with $X_N=-1.3+3Z$.
\QuestionsTasks{
{9}/{sol:q15}/{\textbf{Problem 15.} Find the cutoff $x$ below which $5\%$ of deviations fall. Find both model results, calculate $E_A$, $E_{Pt}$, and $E_{Pn}$. \score{2}}
}
\needspace{5\baselineskip}
\phantomsection
\label{sol:q15}
\subsection*{C9 -- Problem 15: Inverse lower-tail cutoff solution \score{2}}
Here $\nu=9-1=8$; the cumulative target is $0.0500$, with $c=(x+1.3)/3$.
Row $8$, column $\alpha=0.05$, gives $t^*=1.860$, so $c_T=-1.860$.
For normal, row $1.6$, columns $0.04$ and $0.05$, gives $0.9495$ and $0.9505$.
The negative cumulative entries are $0.0505$ and $0.0495$, equally close to $0.0500$.
The tie rule selects $c_N=-1.64$.
\[
 \begin{aligned}
 x_T&=-1.3+3(-1.860)=-6.8800\approx\boxed{-6.88\text{ dollars}},\\
 x_N&=-1.3+3(-1.64)=-6.22=\boxed{-6.22\text{ dollars}}.
 \end{aligned}
\]
Use unrounded contextual results:
\[
 \begin{aligned}
 E_A&=|-6.8800-(-6.22)|=0.6600\approx\boxed{0.66\text{ dollars}},\\
 E_{Pt}&=\frac{100(0.6600)}{|-6.8800|}\approx\boxed{9.59\%},\\
 E_{Pn}&=\frac{100(0.6600)}{|-6.22|}\approx\boxed{10.61\%}.
 \end{aligned}
\]
About $5\%$ of batches save more than $\$6.88$ under $t$, or $\$6.22$ under normal.
The normal threshold is about $\$0.66$ closer to the mean deviation.
\secInicio[sec:q15]
}

\needspace{8\baselineskip}
\sectionScored{Order processing cutoff}{C9-\sml{2}}{sec:q16}
{
\textbf{C9 -- Inverse upper-tail cutoff.} A sample of $n=18$ orders has processing-time deviation from target, in minutes, with mean $m=2.4$ and approximate standard deviation $a=4$. Compare $X_T=2.4+4T$, $T\sim t_{n-1}$, with $X_N=2.4+4Z$.
\QuestionsTasks{
{9}/{sol:q16}/{\textbf{Problem 16.} Find the cutoff $x$ exceeded by $10\%$ of deviations. Find both model results, calculate $E_A$, $E_{Pt}$, and $E_{Pn}$. \score{2}}
}
\needspace{5\baselineskip}
\phantomsection
\label{sol:q16}
\subsection*{C9 -- Problem 16: Inverse upper-tail cutoff solution \score{2}}
Here $\nu=18-1=17$; convert right tail $0.10$ to cumulative target $0.9000$.
With $c=(x-2.4)/4$, row $17$, column $\alpha=0.10$, gives $c_T=1.333$.
For normal, row $1.2$, column $0.08$, gives $\Phi(1.28)=0.8997$, closer to $0.9000$ than $\Phi(1.29)=0.9015$ (column $0.09$).
Thus $c_N=1.28$.
\[
 \begin{aligned}
 x_T&=2.4+4(1.333)=7.7320\approx\boxed{7.73\text{ minutes}},\\
 x_N&=2.4+4(1.28)=7.52=\boxed{7.52\text{ minutes}}.
 \end{aligned}
\]
Use unrounded cutoffs:
\[
 \begin{aligned}
 E_A&=|7.7320-7.52|=0.2120\approx\boxed{0.21\text{ minutes}},\\
 E_{Pt}&=\frac{100(0.2120)}{|7.7320|}\approx\boxed{2.74\%},\\
 E_{Pn}&=\frac{100(0.2120)}{|7.52|}\approx\boxed{2.82\%}.
 \end{aligned}
\]
About $10\%$ of orders take more than $7.73$ minutes above target under $t$, or $7.52$ minutes under normal.
The normal cutoff is about $0.21$ minutes lower.
\secInicio[sec:q16]
}
\needspace{12\baselineskip}
\sectionScored{Research option}{C10-\sml{6}}{sec:q17}
{
\textbf{Exam:}
\QuestionsTasks{
{10}/{sol:q17a}/{Select one model and describe plausible real-life data in 1--2 sentences. \score{6}}
}
\begin{enumerate}[label=\textbf{(\alph*)},leftmargin=*]
\item Linear probability density over a finite interval.
\item Exponential distribution (waiting times).
\item Normal distribution (measurements).
\end{enumerate}

\textbf{Online report:} Submit a PDF through \textbf{Open} using real economic or financial data and the selected model. Include simulation, sampling, portfolio risk, methods, assumptions, limitations, and evidence-based conclusions. The report is the C10 evidence.

\needspace{6\baselineskip}
\phantomsection
\label{sol:q17a}
\subsection*{C10 -- Acceptable alternative brief responses \score{6}}

This selection task has no unique answer. For this worked response, select \textbf{(b)} and use its two-sentence response below.
A student submits \textbf{only one} 1--2 sentence response; (a) and (c) are clearly identified \textbf{alternatives}:
\begin{enumerate}[label=\textbf{(\alph*)},leftmargin=*]
\item \textbf{Linear density.} I would study customers' spending in thousands of Colombian pesos on a promotional service priced continuously between 20 and 60 thousand pesos, proposing a linearly decreasing density over $[20,60]$. I would check whether real observations support that proposed model.
\item \textbf{Exponential waiting times.} I would study the minutes between payment requests reaching a small business during a stable trading period, proposing an exponential waiting-time model if arrivals are independent and occur at a constant average rate. I would check that assumption with real observations.
\item \textbf{Normal measurements.} I would study the fill volumes in millilitres of bottles sold by a beverage producer, proposing a normal model during stable production. I would check whether real measurements support that model.
\end{enumerate}
Each describes possible real-life economic data and a \emph{proposed} model requiring validation;
none claims that uncollected data follow the model. Other plausible responses are acceptable.
Each example is a brief exam response, not an empirical analysis or a completed online report.
The exam prints six points beside the brief selection and identifies the separate online report as the C10 evidence; no separate report point allocation is stated.

\textbf{Separate online PDF report (not the brief exam response).}
Submit through \textbf{Open}, focusing on the \textbf{one selected option}.
Investigate an economic or financial question using real data and the selected continuous model.
Include simulation, sampling analysis, and portfolio-risk measures; justify methods and present results,
assumptions, limitations, and evidence-based conclusions. These are the broader C10 requirements;
no completed report, dataset, simulation, portfolio analysis, or empirical conclusion is supplied by these examples.

\secInicio[sec:q17]
}

\end{document}
